\rtmtight
\begin{tblr}{width=\textwidth,colspec={|Q[c,m,2.1cm]|X[l,m]|},hlines,vlines,hspan=minimal,rowsep=1.5pt,colsep=3pt,row{1}={bg=headgray,font=\bfseries}}
\SetCell{c,m}\hfil\logo\hfil & \textbf{POLITEKNIK NEGERI MALANG}\par\textbf{JURUSAN TEKNOLOGI INFORMASI}\par\textbf{PROGRAM STUDI: D4 TEKNIK INFORMATIKA} \\
\SetCell[c=2]{l} \textbf{RENCANA TUGAS MAHASISWA (RTM) DAN RUBRIK PENILAIAN} & \\
Nama Mata Kuliah & Aljabar Linier \\
Kode Mata Kuliah & RTI252002 \quad \textbf{sks / Jam perminggu:} 2 SKS/4 jam \quad \textbf{Semester:} 2 \\
Dosen Pengampu & \begin{enumerate}\item Muhammad Afif Hendrawan, S.Kom., M.T.\item Prof. Dr. Eng. Cahya Rahmad, S.T., M.T.\item Erfan Rohadi, S.T., M.Eng., Ph.D.\item Mungki Astiningrum, S.T., M.Kom.\item Triana Fatmawati, S.T., M.T.\item Adevian Fairuz Pratama, S.ST., M.Eng.\end{enumerate} \\
\end{tblr}
\assessmentblock{Tugas Worksheet Aljabar Linier}{Tugas terstruktur (worksheet dan latihan)}{SCPMK0902-01101, SCPMK0902-01102, SCPMK0902-01103}{Mahasiswa menyelesaikan worksheet mingguan (Worksheet 1--5) yang mencakup permodelan dan penyelesaian SPL, operasi matriks, invers, determinan, aturan Cramer, vektor, eigen, proyeksi ortogonal, serta studi kasus pemanfaatan aljabar linier di bidang informatika.}{Mengerjakan latihan terstruktur di kelas, menyelesaikan worksheet secara mandiri dengan langkah pengerjaan runtut, memvalidasi hasil secara komputasional, dan mengumpulkan melalui LMS.}{Lembar jawaban worksheet dengan langkah penyelesaian lengkap dan hasil validasi.}{Ketepatan penyelesaian SPL (pemodelan, eliminasi, REF/RREF) & 5 & Model SPL benar, langkah eliminasi/operasi baris elementer runtut, himpunan penyelesaian benar dan tervalidasi. & Model SPL benar tetapi terdapat kesalahan minor pada langkah eliminasi atau validasi hasil belum lengkap. & Model SPL keliru atau langkah eliminasi tidak runtut sehingga penyelesaian salah. \\ \\
Ketepatan operasi matriks, invers, determinan, dan aturan Cramer & 9 & Operasi matriks, inversi, perhitungan determinan dengan ekspansi kofaktor, dan aturan Cramer dikerjakan benar dengan langkah lengkap. & Sebagian besar operasi benar, tetapi ada kesalahan perhitungan minor atau langkah yang tidak ditunjukkan. & Banyak kesalahan konsep atau perhitungan pada operasi matriks, invers, determinan, atau aturan Cramer. \\ \\
Ketepatan operasi vektor, eigen, dan proyeksi ortogonal & 4 & Operasi dot/cross product, perhitungan eigenvalue/eigenvector, dan proyeksi ortogonal benar disertai interpretasi geometri. & Operasi vektor sebagian benar, tetapi interpretasi atau perhitungan eigen/proyeksi masih ada kesalahan minor. & Operasi vektor, eigen, atau proyeksi ortogonal banyak yang salah atau tidak dikerjakan. \\ \\
Penyelesaian studi kasus informatika & 2 & Konsep aljabar linier dipilih tepat, diterapkan pada kasus informatika, dan hasilnya diinterpretasikan dengan benar. & Konsep yang dipilih relevan, tetapi penerapan atau interpretasi hasil kurang lengkap. & Konsep yang dipilih tidak relevan atau penerapan pada kasus salah. \\}

\assessmentblock{Kuis 1 SPL dan Eliminasi}{Kuis (ujian tertulis)}{SCPMK0902-01101}{Mahasiswa mengerjakan kuis tertulis materi pekan 1--3: permodelan SPL, notasi matriks, operasi baris elementer, metode Gauss dan Gauss-Jordan.}{Mengerjakan soal kuis secara mandiri, closed book, sesuai jadwal asesmen.}{Lembar jawaban kuis dengan langkah penyelesaian.}{Ketepatan pemodelan SPL dan notasi matriks & 4 & Permasalahan dimodelkan menjadi SPL dan notasi matriks augmented dengan benar dan lengkap. & Model SPL benar tetapi notasi matriks kurang tepat atau tidak lengkap. & Permasalahan tidak dapat dimodelkan menjadi SPL atau notasi matriks salah. \\ \\
Ketepatan penyelesaian SPL dengan Gauss dan Gauss-Jordan & 6 & Langkah REF/RREF runtut, hasil akhir benar, dan sesuai kunci jawaban. & Langkah eliminasi sebagian besar benar, tetapi terdapat kesalahan perhitungan pada hasil akhir. & Langkah eliminasi salah arah atau hasil akhir tidak sesuai kunci jawaban. \\}

\assessmentblock{Ujian Tengah Semester}{UTS (ujian tertulis)}{SCPMK0902-01101, SCPMK0902-01102}{Mahasiswa mengerjakan ujian tertulis materi pekan 1--7: permodelan dan penyelesaian SPL, metode Gauss/Gauss-Jordan, operasi matriks, inversi, partisi matriks, determinan, dan ekspansi kofaktor.}{Mengerjakan soal ujian secara mandiri, closed book, sesuai jadwal asesmen.}{Lembar jawaban ujian dengan langkah penyelesaian.}{Ketepatan penyelesaian SPL dan eliminasi Gauss/Gauss-Jordan & 12.5 & Model SPL, langkah REF/RREF, dan himpunan penyelesaian benar serta divalidasi. & Langkah eliminasi sebagian besar benar dengan kesalahan perhitungan minor. & Model atau langkah eliminasi salah sehingga penyelesaian tidak benar. \\ \\
Ketepatan operasi matriks, invers, dan determinan & 12.5 & Operasi matriks, inversi $n \times n$, dan determinan dengan ekspansi kofaktor dikerjakan benar dan runtut. & Sebagian operasi benar, tetapi ada kesalahan pada inversi atau perhitungan determinan. & Operasi matriks, invers, atau determinan banyak yang salah atau tidak dikerjakan. \\}

\assessmentblock{Kuis 2 Matriks dan Vektor}{Kuis (ujian tertulis)}{SCPMK0902-01102, SCPMK0902-01103}{Mahasiswa mengerjakan kuis tertulis materi pekan 9--11: aturan Cramer, vektor, dot product, dan cross product.}{Mengerjakan soal kuis secara mandiri, closed book, sesuai jadwal asesmen.}{Lembar jawaban kuis dengan langkah penyelesaian.}{Ketepatan penyelesaian SPL dengan aturan Cramer & 5 & Determinan pembentuk dihitung benar dan solusi SPL dengan aturan Cramer tepat. & Prosedur aturan Cramer benar, tetapi ada kesalahan perhitungan determinan. & Prosedur aturan Cramer tidak dipahami atau solusi salah. \\ \\
Ketepatan operasi vektor dot dan cross product & 5 & Operasi dot/cross product benar dan digunakan tepat untuk menyelesaikan permasalahan vektor. & Operasi vektor sebagian benar, tetapi penerapan pada permasalahan kurang tepat. & Operasi dot/cross product salah atau tidak dikerjakan. \\}

\assessmentblock{Ujian Akhir Semester}{UAS (ujian tertulis)}{SCPMK0902-01101, SCPMK0902-01102, SCPMK0902-01103}{Mahasiswa mengerjakan ujian tertulis materi pekan 1--16 yang mencakup SPL, operasi matriks, determinan, aturan Cramer, vektor, eigenvalue/eigenvector, proyeksi ortogonal, dan penerapannya pada kasus informatika.}{Mengerjakan soal ujian secara mandiri, closed book, sesuai jadwal asesmen.}{Lembar jawaban ujian dengan langkah penyelesaian.}{Ketepatan penyelesaian SPL dan validasi hasil & 10 & SPL dimodelkan dan diselesaikan benar dengan metode yang sesuai serta hasil divalidasi. & Penyelesaian SPL sebagian besar benar dengan kesalahan minor pada validasi. & Penyelesaian SPL salah atau tanpa langkah yang dapat diperiksa. \\ \\
Ketepatan operasi matriks, determinan, dan aturan Cramer & 12.5 & Operasi matriks, invers, determinan, dan aturan Cramer dikerjakan benar dan runtut. & Sebagian besar benar, tetapi terdapat kesalahan perhitungan minor. & Banyak kesalahan konsep atau perhitungan pada matriks dan determinan. \\ \\
Ketepatan vektor, eigen, dan proyeksi ortogonal serta penerapannya & 12.5 & Perhitungan vektor, eigenvalue/eigenvector, dan proyeksi ortogonal benar serta diterapkan tepat pada kasus informatika. & Perhitungan sebagian benar, tetapi penerapan pada kasus kurang tepat. & Perhitungan vektor/eigen/proyeksi salah atau penerapan tidak relevan. \\}

\begin{tblr}{width=\textwidth,colspec={|Q[l,m,3.2cm]|X[l,m]|},hlines,vlines,hspan=minimal,row{1-3}={bg=headgray},rowsep=1.5pt,colsep=3pt}
Jadwal Pelaksanaan: & Minggu 1--17 sesuai RPS. Setiap evaluasi memiliki RTM dan rubrik multi-kriteria. \\
Lain-Lain Yang Diperlukan: & LMS, worksheet, dan perangkat presentasi. \\
Pustaka: & Mengacu pustaka utama dan pendukung pada bagian RPS. \\
\end{tblr}
